\begin{table}[H]
    \centering
    \caption{Analisis Sensitivitas Thorndike Case III terhadap Perlakuan LBI pada Skor Komposit}
    \label{tab:sens_thorndike}
    \begin{spacing}{1.0}
    \resizebox{\textwidth}{!}{%
    \begin{tabular}{lccccccc}
        \toprule
        Skenario & $k$ & $\bar{r}$ \textit{Raw Sample} & $s_Z$ Sampel & $S_Z$ Populasi & $\lambda^2$ & $\bar{R}$ Terkoreksi & Rasio $SD$ (min--maks) \\
        \midrule
        8 subtes (spesifikasi utama) & 8 & 0,300 & 53,84 & 67,72 & 0,582 & 0,425 & 0,674--1,498 \\
        7 subtes tanpa LBI Soshum & 7 & 0,330 & 57,71 & 68,03 & 0,390 & 0,424 & 0,656--1,450 \\
        7 subtes tanpa LBI Saintek & 7 & 0,307 & 55,47 & 67,43 & 0,478 & 0,414 & 0,666--1,475 \\
        6 subtes tanpa kedua LBI & 6 & 0,349 & 60,43 & 68,31 & 0,278 & 0,421 & 0,645--1,419 \\
        \bottomrule
    \end{tabular}%
    }
    \end{spacing}
    \vspace{2mm}
    \begin{flushleft}
    \footnotesize\textit{Catatan.} Skor komposit adalah rerata $k$ subtes yang dipakai, dan tiap subtes memakai $SD$ nasional tidak resminya sendiri. $\bar{r}$ adalah rerata korelasi \textit{raw sample}, dan $\bar{R}$ adalah rerata korelasi terkoreksi \textit{restriction of range}. $\lambda^2$ adalah multiplikator seleksi ($\frac{S_Z^2}{s_Z^2} - 1$). Rasio adalah $SD$ tersirat populasi terhadap $SD$ nasional tidak resmi (lihat Tabel~\ref{tab:diagnostik}). Rerata kedua skor LBI tidak dipakai karena $SD$ nasionalnya tidak tersedia.
    \end{flushleft}
\end{table}
